\documentclass{gescons}

\genre {Entrevista}
\author{Isabela Collares}
\title{A Lógica do Pensene (Ou a~descoberta sobre quase todas as coisas)}


\begin{document}
    \makeentrevistatitle
    \coverart{back/Isabela_Collares}

    \begin{multicols}{3}

\begin{center}
    \includegraphics[width=4cm]{articles/entrevista/mockups/Isabela_Collares.png}
\end{center}
\vspace{-1cm}


\subsubsection{1. Qual foi a~motivação para a~escrita da obra? Por que a~definição deste tema para publicação de um livro?}

\sloppy

A motivação principal foi preencher uma lacuna que eu observava no acervo literário da Ciência: temos uma produção editorial bastante consolidada voltada ao público infantil e~também aos adultos, mas havia uma escassez de obras direcionadas aos adolescentes.

\fussy

Essa fase da vida é~decisiva no processo evolutivo e,~muitas vezes, é~o~momento em que surgem crises, dúvidas e~desafios ligados à~identidade e~às emoções.

A proposta da ``A Lógica do Pensene'' é~justamente oferecer a~esse público uma ferramenta de autocompreensão --- para que o~jovem perceba, desde cedo, a~relação íntima e~lógica entre pensamento, sentimento e~energia.

Compreender o~funcionamento do pensene é~uma senha profilática poderosa, capaz de ajudar o~adolescente a~minimizar a~fase do porão consciencial e~a~lidar de forma mais lúcida com as próprias manifestações.


\subsubsection{2. Quais foram as principais percepções, intra e~extrafísicas, durante a~pesquisa e~a~escrita da obra? E~posterior ao lançamento?}

O lançamento está previsto para dezembro, mas durante o~processo de escrita as percepções foram muito nítidas, tanto no intrafísico quanto no extrafísico.

Extrafisicamente, senti um acoplamento mais ostensivo com o~amparador, como se o~livro realmente tivesse sido escrito ``a muitas mãos''.

Inicialmente, eu me preparava para escrever uma obra sobre pensene voltada às crianças, em parceria com a~colega intermissivista Aline Niemeyer. Porém, em uma noite específica, recebi a~inspiração clara de que o~livro deveria ser voltado para adolescentes --- e~as ideias vieram prontas: os personagens, o~enredo e~até as conexões com Minas Gerais e~com minha própria família.

No intrafísico, a~escrita coincidiu com o~período em que assumi turmas de ensino médio e~atividades docentes na universidade, o~que me permitiu vivenciar o~contato diário com adolescentes e~entender mais de perto suas formas de pensar e~sentir. Tudo isso contribuiu muito para a~coerência da obra.

\subsubsection{3. Qual o~maior aprendizado com a~escrita desta obra?}

O maior aprendizado foi perceber, na prática, como a~escrita conscienciológica reverbera nas várias dimensões da nossa vida.

O livro teve desdobramentos importantes no meu núcleo familiar, especialmente relacionados a~temas e~lugares ligados a~Minas Gerais e~ao Vale do Paraíba --- regiões que o~meu pai, historiador, pesquisou profundamente ao longo da vida.

Sem saber do conteúdo do livro, ele começou a~compartilhar comigo informações e~reflexões que tinham relação direta com a~obra.

Foi uma grande sincronicidade. Isso me mostrou como o~posicionamento autoral reverbera energeticamente e~como a~escrita pode ser um agente de reconciliação e~recomposição grupocármica.

\subsubsection{4. O~que poderia dizer como incentivo para que mais pesquisadores invistam na publicação de obras conscienciológicas?}

A escrita conscienciológica é~uma oportunidade de reciclagem pessoal e~grupal.

Cada autor que se posiciona para escrever ativa um campo de amparo e~favorece encaminhamentos, tanto próprios quanto de familiares e~consciências afins.

Escrever é~um ato de autopesquisa e~de interassistência.

Quando entendemos que a~obra não é~apenas nossa, mas um trabalho a~muitas mãos --- envolvendo consciências intra e~extrafísicas ---, fica claro o~quanto vale a~pena investir nesse processo, deixando um legado evolutivo para as novas gerações e~contribuindo de forma singular para a~série de revezamentos multiexistenciais.

    
    \end{multicols}
\end{document}
